\section{What the stipulated RH and GRH inputs actually discharge}

\paragraph{Input convention.}
Write A for the supplied Alp\"oge theorem: RH for $\zeta$ and GRH for all Dirichlet $L$-functions. The relevant exported declarations are \texttt{riemannHypothesis}, \texttt{LFunction\_eq\_zero\_re}, and \texttt{LFunction\_eq\_zero\_re\_of\_isPrimitive} in \texttt{ALPOGE\_Main.lean}; their predicates are in \texttt{ALPOGE\_Statements.lean}. We use A as stipulated and do not claim to have rebuilt its imported archive. Write M for Masullo's claimed GRH for Hecke $L$-functions over every number field, cited conditionally as the prompt's \S1.2. The supplied index does not contain a numbered \S1.2 statement or the full proof, so no theorem number is invented for M. Theorem~13.3.1.1 (PDF p.~522 in the index) is a trace-formula claim; Conjecture~14.1.2.1 (PDF p.~529) explicitly concerns automorphic $\mathrm{GL}_n$ GRH and simplicity. Neither is silently substituted for M or for a proof of general automorphic GRH.

\subsection{P-Z01: Montgomery's supported pair correlation}
For an even Schwartz function $f$ with $\operatorname{supp}\widehat f\Subset(-1,1)$, set $W(u)=4/(4+u^2)$ and $N_0(T)=T\log T/(2\pi)$. Montgomery's theorem and Fourier identity (1973, \S1, pp.~181--182, equations~(1)--(4); P-B05), combined with A, give
\[
 \lim_{T\to\infty}\frac1{N_0(T)}
 \sum_{0<\gamma,\gamma'\le T}
 f\left(\frac{\log T}{2\pi}(\gamma-\gamma')\right)W(\gamma-\gamma')
 = f(0)+\int_{\mathbb R}f(u)\left[1-\left(\frac{\sin\pi u}{\pi u}\right)^2\right]du.
\]
The sum is over ordered zero indices with multiplicity. Subtracting the identical-index diagonal gives the sine-kernel two-point density. To check the normalization directly, the right side equals
$\widehat f(0)+\int_{-1}^1|u|\widehat f(u)\,du$ under the support restriction, exactly the contribution of Montgomery's delta term at frequency zero and its $|u|$ term. This is now unconditional under the input convention. The unrestricted-support assertion, the complete nearest-neighbour spacing law, and all-height GUE statistics require additional estimates. The missing analytic input is control of the off-diagonal prime correlations occurring beyond the supported range; critical-line location does not provide that estimate in the cited theorem.

\subsection{P-Z02: the Hughes--Rudnick hypothesis correction}
Put $x_{\chi,j}=\gamma_{\chi,j}\log q/(2\pi)$ and $W_f(\chi)=\sum_jf(x_{\chi,j})$. Hughes--Rudnick (2003), Introduction, expressly permits complex $\gamma_{\chi,j}$ in $\rho=1/2+i\gamma_{\chi,j}$ (P-B06). Its analytic mean, variance, and restricted moments are therefore already unconditional before A. For admissible $f$ (Definition~2.1), averaged over the $q-2$ nonprincipal characters modulo a prime $q$, Theorems~3.1, 4.1, and 5.1 respectively give
\[
 \mathbb E W_f\longrightarrow\int f,\qquad
 \operatorname{Var}(W_f)\longrightarrow\int_{-1}^1|u|\widehat f(u)^2\,du,
\]
with supports $[-2,2]$ and $[-1,1]$, and Gaussian centred moments through order $m$ when $\operatorname{supp}\widehat f\subseteq[-\alpha,\alpha]$ and $m\alpha<2$. A makes these ordinates real. It also discharges the explicit GRH in Theorem~8.1, Corollary~8.2 ($\liminf_q\min_{\chi\ne\chi_0}x_{\chi,1}\le1/4$), and Theorem~8.3 (its positive-proportion bound for $\beta\ge0.633$). Here $x_{\chi,1}$ is the first nonnegative scaled ordinate. A finite collection of Gaussian moments is not a Gaussian limit theorem at microscopic scale: the same paper calls the phenomenon mock-Gaussian, and the full random-matrix moments are not all Gaussian.

\subsection{Dependency table}
\begingroup\small
\begin{longtable}{p{0.25\textwidth}p{0.33\textwidth}p{0.34\textwidth}}
\hline
Statement and locator & Exact RH/GRH input or other gate & Status under A and M\\\hline
\endfirsthead
\hline
Statement and locator & Exact RH/GRH input or other gate & Status under A and M\\\hline
\endhead
Montgomery (1973), \S1 theorem and (3)--(4), supported pair correlation & RH for $\zeta$; unchanged $(-1,1)$ Fourier support & Discharged by A (P-Z01).\\
Rudnick--Sarnak (1996), Theorem~1.2, $m=1$ & RH for the fixed primitive Dirichlet $L$-function (or $\zeta$); support $\sum|\xi_j|<2$ & Discharged by A. Theorem~1.1 already proves a smoothed analytic version with complex ordinates.\\
Hughes--Rudnick (2003), Theorems~3.1,4.1,5.1 & No GRH for the analytic statistic; A makes all ordinates real & Already unconditional formulas; real-zero interpretation now supplied by A (P-Z02).\\
Hughes--Rudnick (2003), \S8 & Dirichlet GRH & Discharged by A; original support and constants retained.\\
Gao--Zhao (2020), Theorems~1.1 and~1.3 & GRH for the quadratic/quartic Hecke families over $\mathbb Q(i)$ and auxiliary Hecke characters in Lemma~2.8 & Conditional on M; symplectic density for quadratic family with support $(-2,2)$, unitary for quartic with $(-20/19,20/19)$. A alone does not cover these general Hecke characters.\\
Iwaniec--Luo--Sarnak (2000), Theorems~1.1--1.3 & Standing RH for target cuspidal $\mathrm{GL}_2$ functions, plus Dirichlet RH for extended support; see p.~56, and pp.~88--89 for auxiliary symmetric-square estimates & Beyond A and M for the general family: missing RH for $L(s,f)$. Extra symmetric-square/Rankin--Selberg hypotheses used in individual estimates can be avoided by the averaged replacements explicitly noted in \S4.\\
Iwaniec--Luo--Sarnak (2000), Theorems~1.4--1.5 & RH for target $L(s,\operatorname{sym}^2f)$ ($\mathrm{GL}_3$); auxiliary hypotheses as specified in \S4 & Beyond A and M; missing target symmetric-square RH.\\
Iwaniec--Luo--Sarnak (2000), Theorem~1.9 & Weighted, parity-unsplit analytic one-level statistic, support $(-2,2)$ & Already unconditional; target zero ordinates may be complex.\\
Iwaniec--Luo--Sarnak (2000), Theorem~1.10 & Hypothesis S, a prime exponential-sum estimate & Beyond A and M; the precise missing input is Hypothesis S, not merely another GRH label.\\
Katz--Sarnak, Theorem~9.2.6; Weil II, Theorem~3.5.3 & Purity, the actual geometric/arithmetic monodromy hypotheses, and the stated field-extension limit & Already unconditional in that scope. Growing rank/genus/conductor requires separate uniformity.\\
Sawin (2020), Hypothesis~1.3, Theorem~1.4 (label \texttt{main}), Remark~1.9(1),(5) & $H(n,r_+,r_-,n-1)$ follows from duality and gives fixed-parameter $O(q^{-1/2})$; stronger degree-aspect conclusions require vanishing and stronger uniform Betti bounds & Basic $q$-aspect comparison already unconditional. General stronger cohomology and uniformity gates remain; the unconditional unbalanced case is recorded below.\\
Keating--Snaith / CFKRS general number-field moment predictions & Arithmetic moment asymptotics and off-diagonal terms; CFKRS labels its zeta assertion a conjecture & A and M do not discharge the moment recipe. The finite-$N$ CUE integrals are separate exact theorems.\\
Berry--Keating $H=xp$ and BGS universality heuristics & A self-adjoint spectral realization with the required boundary/domain conditions and appropriate chaotic universality theorem & Beyond A and M; RH alone is not that construction.\\\hline
\end{longtable}
\endgroup

Sawin also proves an unconditional unbalanced case: Corollary~1.8 (repeated at source label \texttt{combined-result}) treats $r_-=1$ and $r_+=r$, with $n>2r+1$, the same stated characteristic exclusions, and error
\[
 O\!\left(\sqrt q\,\bigl(q^{-(p-2r)/(2pr)}(r+2)^{r+1}\bigr)^n n^{r+1}\right),
 \qquad p=\operatorname{char}\mathbb F_q.
\]
For fixed $r,p,q$ satisfying $q^{(p-2r)/(2pr)}>(r+2)^{r+1}$, this error decays exponentially in $n$ up to the displayed polynomial factor. Thus some degree-aspect cases are already proved; the unresolved general moment claim is broader than this restricted theorem.

For Rudnick--Sarnak in degree $m>1$, Theorem~1.2 also needs the target RH and, for $m>3$ in the stated source, their coefficient Hypothesis~(1.7). M covers a Hecke case only after all the relevant identification and other hypotheses have been verified; it is not a blanket automorphic GRH theorem. For ILS, ``Hecke cusp form'' refers to a $\mathrm{GL}_2$ automorphic object; it must not be confused with a one-dimensional Hecke character over a number field.

\paragraph{Quantum-chaos source discipline.}
The connection recalled from memory is the symmetry-class heuristic: chaotic quantum systems without an antiunitary time-reversal symmetry suggest unitary statistics, while spinless time-reversal invariance with $T^2=+1$ suggests orthogonal statistics (the $T^2=-1$ symmetry class is symplectic). The original BGS paper (Bohigas, Giannoni, \& Schmit, 1984, pp.~1--4) concerns chaotic quantum spectra; it does not construct the zeta-zero operator. The bibliographic locator and Berry--Keating's explicit $H_{\rm cl}=XP$ speculation were checked against the primary publication records. Berry and Keating (1999), \emph{The Riemann zeros and eigenvalue asymptotics}, SIAM Review 41, 236--266, and their \emph{$H=xp$ and the Riemann zeros} contribution in \emph{Supersymmetry and Trace Formulae: Chaos and Disorder}, pp.~355--367, are the relevant sources. Detailed internal section/equation locators for the latter were not inspected here, so they are not fabricated.

There is a distinction from bare operator existence. Once real zero ordinates $\gamma_j$ are given, the diagonal operator $He_j=\gamma_je_j$ on $\ell^2$, with domain $\{(a_j):\sum_j\gamma_j^2|a_j|^2<\infty\}$, is self-adjoint. Indeed testing the adjoint against the coordinate vectors forces the same diagonal action and domain. This merely packages the supplied spectrum. It does not supply a specified differential or pseudodifferential Hamiltonian, arithmetic trace formula, boundary conditions, or a proof of chaotic local statistics. Those are the substantive missing constructions in the Berry--Keating/Hilbert--P\'olya direction.
